\documentclass[12pt, a4paper]{article}

% ── Codificación y lengua ─────────────────────────────────────────────────────
\usepackage[utf8]{inputenc}
\usepackage[T1]{fontenc}
\usepackage[spanish]{babel}

% ── Tipografía y diseño ───────────────────────────────────────────────────────
\usepackage{lmodern}
\usepackage{microtype}
\usepackage{geometry}
\geometry{
  a4paper,
  top=2.5cm, bottom=2.5cm,
  left=3cm,  right=2.5cm
}
\usepackage{setspace}
\onehalfspacing

% ── Sangría en todos los párrafos (incluido el primero de cada sección) ───────
\usepackage{indentfirst}
\setlength{\parindent}{1.5em}
\setlength{\parskip}{4pt}

% ── Colores (solo para elementos decorativos, no para texto) ──────────────────
\usepackage[dvipsnames]{xcolor}
\definecolor{azuloscuro}{RGB}{31, 56, 100}
\definecolor{azulmedio}{RGB}{46, 117, 182}
\definecolor{grisclaro}{RGB}{245, 245, 245}
\definecolor{amarillotodo}{RGB}{255, 243, 205}

% ── Encabezados y secciones (texto en negro) ──────────────────────────────────
\usepackage{titlesec}
\titleformat{\section}
  {\large\bfseries}{{\thesection.}}{0.6em}{}
  [\vspace{-4pt}\color{azulmedio}\rule{\linewidth}{1pt}\vspace{2pt}]
\titleformat{\subsection}
  {\normalsize\bfseries}{{\thesubsection.}}{0.5em}{}
\titleformat{\subsubsection}
  {\normalsize\itshape}{{\thesubsubsection.}}{0.5em}{}
\titlespacing{\section}{0pt}{18pt}{8pt}
\titlespacing{\subsection}{0pt}{12pt}{4pt}

% ── Encabezado y pie de página ────────────────────────────────────────────────
\usepackage{fancyhdr}
\pagestyle{fancy}
\fancyhf{}
\fancyhead[L]{\small\textit{Clasificación Automática de Retinopatía Diabética}}
\fancyhead[R]{\small\textit{PIM --- Curso 2025/2026}}
\fancyfoot[C]{\small\thepage}
\renewcommand{\headrulewidth}{0.4pt}
\renewcommand{\footrulewidth}{0pt}
\setlength{\headheight}{14.5pt}

% ── Tablas ────────────────────────────────────────────────────────────────────
\usepackage{booktabs}
\usepackage{array}
\usepackage{tabularx}
\usepackage{multirow}
\usepackage{colortbl}

% ── Figuras ───────────────────────────────────────────────────────────────────
\usepackage{graphicx}
\usepackage{float}
\usepackage{caption}
\captionsetup{
  font=small,
  labelfont=bf,
  labelsep=period
}

% ── Matemáticas ───────────────────────────────────────────────────────────────
\usepackage{amsmath}
\usepackage{amssymb}

% ── Listas ────────────────────────────────────────────────────────────────────
\usepackage{enumitem}
\setlist[itemize]{itemsep=2pt, topsep=4pt}
\setlist[enumerate]{itemsep=2pt, topsep=4pt}

% ── Hipervínculos y referencias ───────────────────────────────────────────────
\usepackage[
  colorlinks=true,
  linkcolor=black,
  citecolor=black,
  urlcolor=black,
  pdftitle={Clasificación Automática de Retinopatía Diabética},
  pdfauthor={García Fernández, Véliz Ayala, Solano Martínez},
]{hyperref}
\usepackage[style=ieee, sorting=none]{biblatex}
\addbibresource{referencias.bib}

% ── Otros ─────────────────────────────────────────────────────────────────────
\usepackage{tocbibind}
\usepackage{xspace}

% ── Comandos propios ──────────────────────────────────────────────────────────
\newcommand{\kappados}{$\kappa^2$\xspace}

% Caja amarilla para marcar contenido pendiente (texto en negro)
\newcommand{\pendiente}[1]{%
  \vspace{4pt}%
  \colorbox{amarillotodo}{%
    \parbox{\dimexpr\linewidth-2\fboxsep}{%
      \small\itshape%
      $\blacktriangleright$ \textbf{[PENDIENTE]:} #1%
    }%
  }%
  \vspace{4pt}%
}

% ══════════════════════════════════════════════════════════════════════════════
\begin{document}
\pagenumbering{gobble}

% ══════════════════════════════════════════════════════════════════════════════
%  PORTADA
% ══════════════════════════════════════════════════════════════════════════════
\begin{titlepage}
  \centering

  \begin{minipage}{\textwidth}
    \color{azuloscuro}\rule{\textwidth}{4pt}
  \end{minipage}

  \vspace{2.5cm}

  {\Huge\bfseries
    Clasificación Automática de\\[6pt]
    Retinopatía Diabética
  }

  \vspace{0.5cm}
  {\large\itshape
    Pipeline Completo de Clasificación por Gradación de Severidad
  }

  \vspace{0.4cm}
  \color{azulmedio}\rule{0.6\textwidth}{1.5pt}\color{black}

  \vspace{2cm}

  \begin{tabular}{rl}
    \textbf{Asignatura:}       & Procesamiento y Análisis de Imagen Médica (PIM) \\[4pt]
    \textbf{Curso académico:}  & 2025 -- 2026 \\[4pt]
    \textbf{Fecha de entrega:} & 24 de abril de 2026 \\
  \end{tabular}

  \vspace{2cm}

  {\large\bfseries Integrantes del grupo}\\[0.6cm]
  \begin{tabular}{c}
    Javier García Fernández    \\[4pt]
    Miguel Ángel Véliz Ayala  \\[4pt]
    Martín Solano Martínez     \\
  \end{tabular}

  \vfill

  \begin{minipage}{\textwidth}
    \color{azuloscuro}\rule{\textwidth}{4pt}
  \end{minipage}

\end{titlepage}

% ══════════════════════════════════════════════════════════════════════════════
%  RESUMEN
% ══════════════════════════════════════════════════════════════════════════════
\newpage
\pagenumbering{roman}

\phantomsection
\section*{Resumen}
\addcontentsline{toc}{section}{Resumen}

Este trabajo aborda la clasificación automática de la retinopatía diabética (RD) en
cinco grados de severidad clínica ---desde grado~0 (sin RD) hasta grado~4
(proliferativa)--- a partir del dataset público APTOS~2019. Se construye un pipeline
completo basado en EfficientNetBN-B3 (MONAI) con preprocesamiento específico para
retinografías mediante máscara circular y CLAHE, y se diseñan y comparan diez
experimentos que abordan las principales limitaciones del modelo base: desbalance de
clases~(Exp.~A), radiómica cuantitativa con PyRadiomics~(Exp.~B), clasificación
zero-shot con BiomedCLIP~(Exp.~C), pipelines jerárquicos~(Exp.~D y E), regresión
ordinal a 384~px con optimización de umbrales~(Exp.~F), corrección de
\textit{shortcut learning} mediante preprocesado Ben~Graham~(Exp.~G), fusión
multimodal CNN\,+\,radiómica con LightGBM~(Exp.~I) y preentrenamiento en
APTOS~2015~(Exp.~J). La métrica principal es el kappa cuadrático ponderado~($\kappa^2$),
estándar en la literatura de gradación automática de RD. Esta métrica no solo mide acierto, sino que penaliza más los errores graves que los leves, teniendo en cuenta el orden entre clases. Los resultados completos y el análisis comparativo de todos los experimentos se presentan en la Sección~4.

% ══════════════════════════════════════════════════════════════════════════════
%  ÍNDICE
% ══════════════════════════════════════════════════════════════════════════════
\newpage
\phantomsection
\tableofcontents

% ══════════════════════════════════════════════════════════════════════════════
%  CUERPO DEL DOCUMENTO
% ══════════════════════════════════════════════════════════════════════════════
\newpage
\pagenumbering{arabic}

% ══════════════════════════════════════════════════════════════════════════════
%  1. INTRODUCCIÓN
% ══════════════════════════════════════════════════════════════════════════════
\section{Introducción}

La \textbf{retinopatía diabética} (RD) es una complicación microvascular de la
diabetes mellitus y una de las principales causas de ceguera prevenible en el mundo.
Con más de 540 millones de adultos diabéticos estimados en 2023 y una prevalencia
de afectación retiniana en torno al 35\,\% de los pacientes~\cite{idf2023,who2023},
el cribado masivo mediante retinografías de fondo de ojo es una necesidad clínica
de primer orden. Sin embargo, la revisión manual de miles de imágenes por
oftalmólogos resulta inviable a gran escala, lo que motiva el desarrollo de sistemas
automáticos de gradación. Además, estos sistemas no solo permiten escalar el
cribado, sino que también facilitan y agilizan el trabajo del personal sanitario,
sirviendo como herramienta de apoyo en la clasificación de la enfermedad y en la
toma de decisiones clínicas. En este contexto, se emplea el dataset público APTOS~2019, ampliamente utilizado en la literatura, que contiene retinografías etiquetadas según el grado de severidad de la enfermedad, lo que permite evaluar de forma objetiva el rendimiento de los modelos propuestos.

La enfermedad progresa en cinco estadios (Tabla~\ref{tab:grados}): desde la ausencia
de lesiones (G0) hasta la neovascularización proliferativa con riesgo de ceguera súbita
(G4), pasando por microaneurismas aislados (G1), hemorragias y exudados duros (G2) y
afectación vascular extensa en los cuatro cuadrantes (G3). Esta naturaleza ordinal del
problema, junto con el severo desbalance de clases del dataset ---G0 representa
$\sim$49\,\% del entrenamiento frente al $\sim$5\,\% de G3---, la ambigüedad entre
grados adyacentes, la pérdida de detalles finos al comprimir las imágenes a baja
resolución y el riesgo de \textit{shortcut learning} sobre artefactos de adquisición,
conforman los desafíos técnicos que este proyecto aborda de forma sistemática.

\begin{table}[H]
  \centering
  \caption{Escala de severidad de la retinopatía diabética.}
  \label{tab:grados}
  \begin{tabularx}{0.9\textwidth}{c l X}
    \toprule
    \rowcolor{azuloscuro!15}
    \textbf{Grado} & \textbf{Nombre} & \textbf{Características clínicas} \\
    \midrule
    0 & Sin RD        & Fondo de ojo normal. Sin lesiones visibles. \\
    \rowcolor{grisclaro}
    1 & Leve          & Microaneurismas aislados. \\
    2 & Moderada      & Hemorragias, exudados duros y anomalías vasculares. \\
    \rowcolor{grisclaro}
    3 & Severa        & Múltiples hemorragias, \textit{venous beading} e IRMA
                        en los cuatro cuadrantes. \\
    4 & Proliferativa & Neovascularización con riesgo de hemorragia vítrea
                        y desprendimiento de retina. \\
    \bottomrule
  \end{tabularx}
\end{table}

\newpage

El objetivo del proyecto es construir un pipeline reproducible y clínicamente
interpretable que parta de un baseline sólido (EfficientNetBN-B3, APTOS~2019,
224~px) y lo mejore de forma justificada mediante diez experimentos que atacan
cada una de las debilidades identificadas. Para ello, se emplea el dataset público
APTOS~2019, compuesto por retinografías etiquetadas en cinco grados de severidad,
lo que permite abordar el problema en un entorno realista y con datos clínicamente
relevantes. La métrica principal es el
\textbf{kappa cuadrático ponderado}~($\kappa^2$), que penaliza los errores
proporcionalmente al cuadrado de la distancia ordinal entre grados, alineándose
con la relevancia clínica de cada tipo de error. La elección de la RD como
problema de estudio responde, además, a la disponibilidad de datos públicos de
calidad (licencia CC~BY~4.0) y a que su naturaleza permite aplicar directamente
las técnicas del curso: preprocesamiento de imagen médica, aprendizaje profundo
con MONAI, radiómica cuantitativa e interpretabilidad con Grad-CAM.

\begin{figure}[H]
  \centering
  \includegraphics[width=\textwidth]{figura1.png}
  \caption{Retinografías representativas por grado de severidad (G0--G4).
    Las diferencias entre grados adyacentes son sutiles, lo que dificulta
    la gradación manual a gran escala.}
  \label{fig:grados}
\end{figure}

El resto del documento se organiza del siguiente modo: la Sección~\ref{sec:arte}
revisa el estado del arte; la Sección~\ref{sec:metodologia} describe los objetivos,
el dataset y el pipeline; la Sección~\ref{sec:resultados} presenta los resultados;
y la Sección~\ref{sec:conclusiones} recoge las conclusiones y vías futuras.

% ══════════════════════════════════════════════════════════════════════════════
%  2. ESTADO DEL ARTE
% ══════════════════════════════════════════════════════════════════════════════
\section{Estado del arte}
\label{sec:arte}

\pendiente{Sección a cargo de Miguel Ángel y/o Martín.\\
  Estructura sugerida: (2.1) métodos clásicos basados en segmentación y radiómica,
  (2.2) redes convolucionales profundas para gradación de RD, (2.3) modelos
  visión-lenguaje y zero-shot, (2.4) radiómica y fusión multimodal.\\
  Referencias clave: Gulshan et al.\ (2016), Tan \& Le EfficientNet (2019),
  BiomedCLIP Zhang et al.\ (2023), Ting et al.\ (2017).}

% ══════════════════════════════════════════════════════════════════════════════
%  3. OBJETIVOS Y METODOLOGÍA
% ══════════════════════════════════════════════════════════════════════════════
\section{Objetivos y metodología}
\label{sec:metodologia}

\pendiente{Sección a cargo del grupo.\\
  Incluir: (3.1) objetivos concretos y medibles, (3.2) descripción del dataset
  APTOS~2019 y APTOS~2015 (fuente, tamaño, modalidad, licencia, distribución de
  clases), (3.3) diagrama del pipeline con justificación de cada decisión de diseño
  (máscara circular, CLAHE, Ben~Graham, EfficientNetBN-B3, regresión ordinal, etc.),
  (3.4) tabla resumen de los diez experimentos.}

% ══════════════════════════════════════════════════════════════════════════════
%  4. RESULTADOS
% ══════════════════════════════════════════════════════════════════════════════
\section{Resultados}
\label{sec:resultados}

\pendiente{Rellenar la tabla y el análisis con los valores de $\kappa^2$ y F1 macro
  obtenidos en la ejecución del notebook (Baseline, Exp.~A--G, I y J).}

% ══════════════════════════════════════════════════════════════════════════════
%  5. CONCLUSIONES Y VÍAS FUTURAS
% ══════════════════════════════════════════════════════════════════════════════
\section{Conclusiones y vías futuras}
\label{sec:conclusiones}

\pendiente{Sección a cargo del grupo.\\
  Incluir: (5.1) hallazgos principales (qué experimento gana y por qué),
  (5.2) limitaciones honestas del trabajo, (5.3) mejoras concretas y líneas futuras.}

% ══════════════════════════════════════════════════════════════════════════════
%  BIBLIOGRAFÍA
% ══════════════════════════════════════════════════════════════════════════════
\newpage
\printbibliography[heading=bibintoc, title={Bibliografía}]

% ══════════════════════════════════════════════════════════════════════════════
%  ANEXO
% ══════════════════════════════════════════════════════════════════════════════
\newpage
\appendix
\section{Tabla de contribución individual}
\label{sec:contribucion}

\begin{table}[H]
  \centering
  \caption{Distribución de tareas entre los miembros del grupo.
    \textbf{R}\,=\,responsable principal, \textbf{P}\,=\,participó.}
  \label{tab:contribucion}
  \begin{tabularx}{\textwidth}{X c c c}
    \toprule
    \rowcolor{azuloscuro!15}
    \textbf{Tarea} &
    \textbf{\small Javier G.} &
    \textbf{\small Miguel Á. V.} &
    \textbf{\small Martín S.} \\
    \midrule
    Exploración del dataset (EDA)           & & & \\
    \rowcolor{grisclaro}
    Preprocesamiento (máscara, CLAHE)       & & & \\
    Dataset class y aumentación             & & & \\
    \rowcolor{grisclaro}
    Baseline (EfficientNetBN-B3)            & & & \\
    Exp. A --- Sobremuestreo                & & & \\
    \rowcolor{grisclaro}
    Exp. B --- Radiómica + LR L1            & & & \\
    Exp. C --- BiomedCLIP zero-shot         & & & \\
    \rowcolor{grisclaro}
    Exp. D --- Pipeline dos etapas          & & & \\
    Exp. E --- Jerárquico tres niveles      & & & \\
    \rowcolor{grisclaro}
    Exp. F --- 384\,px + Regresión ordinal  & & & \\
    Exp. G --- Ben Graham + Hard Mask       & & & \\
    \rowcolor{grisclaro}
    Interpretabilidad Grad-CAM             & & & \\
    Exp. I --- Híbrido CNN + Radiómica      & & & \\
    \rowcolor{grisclaro}
    Exp. J --- Pretrain APTOS 2015          & & & \\
    Redacción de la memoria                 & & & \\
    \rowcolor{grisclaro}
    Preparación de la defensa               & & & \\
    \bottomrule
  \end{tabularx}
\end{table}

\end{document}